\section{Positioning Against Related Work}
\label{sec:disc-related-work}

\Cref{ch:key-concepts-related-work} surveyed the academic and commercial landscape to find the gap this thesis addresses.
Now that the platform is built and measured, we return to that landscape and show where the result sits.
We use the same four threads as the survey: the prior project line, the general-purpose experiment software, the eye-movement analysis literature, and the work on intervention design.
For each one we state plainly what we did, what the prior work did, and why the two differ.

The clearest comparison is with the earlier projects in Reading the Reader, because this thesis inherits their problem rather than a new one.
Each of those projects delivered a capable adaptive reader, but each was extended by copying the previous frontend and changing it, so a new study built on a fork of the last project rather than on a shared contract \cite{refsgaard2024rtr, pereira2024typography, palinec2024reactive}.
Our result keeps the same capability but changes the structure underneath.
Sensing, analysis, decision, and intervention now sit behind build-enforced module boundaries (\cref{sec:eval-modularity}), so the next intervention or decision provider extends the platform in place instead of forking it.
This is the move from a line of one-off prototypes to a platform that builds up over time, which is the precondition the gap analysis named for the systematic comparison the project needs (\cref{subsec:gap-experimentability}).

A researcher who builds a reading study today is more likely to use a general-purpose experiment platform than a prior thesis prototype, so these tools are the harder comparison.
Tobii Pro Lab, PsychoPy, and Psychtoolbox are mature and widely used, and they are very good at designing experiments, presenting stimuli, and capturing gaze data.
But they were not built to adapt the text in response to the reader's gaze, and none of them offers a pluggable decision boundary \cite{tobii2025prolab, peirce2019psychopy, brainard1997psychtoolbox}.
Our platform differs on exactly these two points: it closes the loop from gaze to a typographic change in real time, and it lets an external decision provider connect behind a stable contract.
This is the upper-right region that the capability comparison and the market quadrant of \cref{ch:key-concepts-related-work} leave empty (\cref{tab:capability-matrix}, \cref{fig:market-quadrant}).
We do not overstate the contrast.
Those platforms are general tools, proven by years of use across many kinds of study, whereas ours is a focused tool for one family of gaze-driven reading experiments, so we present it as filling the gap they leave open rather than as a replacement for them.

The eye-tracking literature this thesis builds on treats the reduction of raw gaze to oculomotor events as the first step.
That step uses an explicit identification algorithm of the kind catalogued by Salvucci and Goldberg, read together with the reading-process findings that Rayner consolidated \cite{salvucci2000identifying, rayner1998eyemovements}.
A common theme in that work is that any later result, such as predicting reading performance or telling reading from non-reading, is only trustworthy if this extraction step is fixed and documented, so the result can be compared or replayed \cite{guan2022predictive, landsmann2019classification}.
Our contribution here is architectural rather than algorithmic.
We make the extraction an explicit, replaceable analysis strategy whose events are logged and can be replayed from the session record (\cref{sec:eval-control}, \cref{sec:impl-sensing}), so the step the literature treats as a methodological precondition becomes an inspectable, swappable part of the running system.
We do not claim a better detector: the built-in strategy is a dwell-based, area-of-interest method suited to word-level reading, and we present it as one interchangeable implementation behind the analysis contract, not as an advance in fixation detection (\cref{sec:impl-sensing}).

The last thread is the work on intervention design in and around the project, which asks which typographic changes help or hinder reading and how a reader recovers from them.
Pereira and Andrei compared four intervention forms and reported that they differ a great deal in disruption and acceptance, and Jensen et al. studied how preserving context across a change affects recovery, and found that intervention design measurably shapes the reading experience \cite{pereira2024typography, jensen2025context}.
The typographic parameters these interventions change are themselves grounded in controlled readability studies \cite{rello2016makeitbig, beymer2008eyetracking}.
Our work comes at the same concern from the other side.
Instead of measuring the effect of one intervention, we make an intervention's disruption something the platform measures and records for every change, applying it at a controlled boundary and writing a graded recovery outcome into the session record (\cref{sec:eval-context}, \cref{sec:contributions}).
This complements that line of work rather than competing with it, because the efficacy findings it produces are exactly the kind of result our platform is built to generate in a systematic and reproducible way, which the prior prototypes could not support.

Across the four threads a single pattern appears.
Prior work delivered real capability or real findings, but each one within a fixed structure: a forked prototype, a general-purpose tool that does not adapt the text, an analysis step redone for each study, or a single intervention measured on its own.
What this thesis adds is not another capability of that kind, but the modular structure that makes such capabilities interchangeable and such findings reproducible.
It is that structure, seen as design knowledge rather than as one particular system, that the next section examines.
